%% Cover letter accompanying the MBE submission of paper 2.
%% Built by build.sh alongside the manuscript.
%%
%% BEFORE SENDING, check the "Related manuscripts" paragraph.  It discloses the
%% NeurIPS workshop version, the bioRxiv preprint and the companion MBE
%% submission.  There is also a PSB version of this paper in ../../psb-paper;
%% if that was submitted anywhere, MBE requires it disclosed here too.

\documentclass[11pt]{article}
\usepackage[letterpaper,top=16mm,bottom=14mm,left=22mm,right=22mm]{geometry}
\usepackage[T1]{fontenc}
\usepackage{microtype}
\usepackage[hidelinks]{hyperref}
\pagestyle{empty}
\setlength{\parskip}{0.38em}
\setlength{\parindent}{0pt}

\begin{document}

{\small
Department of Bioengineering\\
College of Engineering\\
University of California, Berkeley\\
374C Stanley Hall, Berkeley, CA 94720, USA\\
\texttt{ihh@berkeley.edu}
}

\hfill\today

Dear Editors,

Please consider the enclosed manuscript, ``Variational inference in coupled
models of amino acid substitution'', for publication in \emph{Molecular Biology
and Evolution} as a Methods paper.

A model of amino acid coevolution has to solve two problems at once: how to
parameterize an interaction between sites without an explosion of free
parameters, and how to compute a likelihood on a state space that grows
exponentially in the number of coupled sites. We address both. Working out
expectation-maximization for coupled continuous-time Markov chains under the
symmetries such a model should obey, reversibility and site exchangeability, we
find a Klein four-group structure on the flux tensor that halves the parameter
count; the resulting pair model fits held-out structural alignment data
slightly better than the current state of the art in coupled amino acid
substitution models, the CherryML Q2 matrix of Prillo et al., with half as many
parameters and similar training data. Fitting one pooled matrix, which is standard practice,
turns out to hide most of the signal. Hydropathy and volume couplings are
strong within individual interaction classes but flip sign between classes and
cancel on pooling; charge is the weaker axis, but it is complementary in almost
every class, so it is the one that survives. A mixture over interaction classes
recovers what pooling destroys. For clusters larger than a pair we give an
evidence lower bound built from the same sufficient statistics, in which the
coupling contracts onto single-site bridge marginals, so the cost is polynomial
rather than exponential in cluster size. It is closed form for pairs and
computed by Gauss-Legendre quadrature otherwise. Against the Euler-Lagrange
mean field of Cohn et al. (2010) it is about a hundred times faster and much
better conditioned; their ODE system failed to converge on 30 to 40 percent of
the pairs we tried at short branch lengths.

Epistasis between contacting residues is well documented, but phylogenetic
practice still assumes sites evolve independently, largely because coupled
models are expensive to fit and to score. The reduction reported here makes
coupled substitution matrices cheap enough to fit on ordinary structural
alignment corpora, and the bound makes them cheap enough to score on clusters
of more than two sites. The finding that pooled fitting cancels real couplings
also bears on how coevolutionary signal is measured in the first place.

Related manuscripts: an abbreviated version of this manuscript has been
submitted for presentation at the AI for Stochastic Dynamics workshop, affiliated
with the NeurIPS meeting in December 2026; that workshop will not publish
proceedings. A companion
manuscript, ``Closed-form Baum--Welch for TKF-based evolutionary models'', is
being submitted to \emph{Molecular Biology and Evolution} at the same time. It
concerns the insertion-deletion side of the same modeling programme and its
results do not overlap with these.

As stated in the Acknowledgements, we used Anthropic's Claude and OpenAI's
ChatGPT in the course of this work, to check mathematical results, write
specifications, develop and run code, and draft text. We checked and edited
all of it, and we are responsible for the content.

The work is original, has not been published elsewhere, and is not under
consideration by another journal. All authors have read and approved the final
version of the manuscript. I am the corresponding author
(\texttt{ihh@berkeley.edu}). We declare no conflict of interest.

Sincerely,

\vspace{1em}

Ian Holmes, B.A., M.A., Ph.D.\\
Professor, Department of Bioengineering\\
University of California, Berkeley

\end{document}
